% ============ 第2章 §2.11 H^p (0<p<1)与积分表示(原书 pp.153-154) ============
\section{\texorpdfstring{$H^p$($0<p<1$)与积分表示}{Hp(0<p<1)与积分表示}}\label{sec:2.11}
\index{Fourier 系数}


$(1+z)(1-z)^{-1}$ 属于 $\bigcap_{p<1} H^{p}$ 这一事实很容易证明.定理~\ref{thm:2.11.2} 利用引理~\ref{lem:2.11.1} 推广了这一结果.

\begin{lemma}\label{lem:2.11.1}
设 $G(z)$ 在 $D$ 上实值且次调和.若 $\phi \in H^{\infty}$,$\|\phi\|_{\infty} \leq 1$ 且 $\phi(0) = 0$,则
\begin{equation*}
\int_{-\pi}^{\pi} G(\phi(re^{i\theta}))\dd \theta \leq \int_{-\pi}^{\pi} G(re^{i\theta})\dd \theta \quad (0 < r < 1).
\end{equation*}
\end{lemma}

\begin{proof}
设 $U(z)$ 在 $|z| < r$ 上调和且在 $|z| = r$ 上 $U(z) = G(z)$.因 $G(z)$ 次调和,在 $|z| \leq r$ 上有 $G(z) \leq U(z)$.由 Schwarz 引理,$|\phi(z)| \leq |z|$,故 $G(\phi(z)) \leq U(\phi(z))$,因此
\begin{align*}
\frac{1}{2\pi}\int_{-\pi}^{\pi} G(\phi(re^{i\theta}))\dd \theta & \leq \frac{1}{2\pi}\int_{-\pi}^{\pi} U(\phi(re^{i\theta}))\dd \theta \\
& = U(\phi(0)) = U(0) = \frac{1}{2\pi}\int_{-\pi}^{\pi} U(re^{i\theta})\dd \theta \\
& = \frac{1}{2\pi}\int_{-\pi}^{\pi} G(re^{i\theta})\dd \theta .
\end{align*}
\end{proof}

\begin{theorem}\label{thm:2.11.2}
若 $f \in H(D)$ 且 $\operatorname{Re} f(z) > 0$($z \in D$),则 $f \in \bigcap_{p<1} H^{p}$.
\end{theorem}

\begin{proof}
因 $\operatorname{Re} f(z) > 0$($z \in D$),可设 $f(0) = 1$.令 $\phi(z) = (f(z)-1)/(f(z)+1)$,则 $\phi \in H^{\infty}$,$\|\phi\|_{\infty} \leq 1$ 且 $\phi(0) = 0$.$f(z) = (1+\phi(z))/(1-\phi(z))$,令
\begin{equation*}
G(z) = \left| \frac{1+z}{1-z} \right|^{p},
\end{equation*}
则 $G(z)$ 次调和且 $|f(z)|^{p} = G(\phi(z))$,故由引理~\ref{lem:2.11.1} 可证 $f \in \bigcap_{p<1} H^{p}$.
\end{proof}

设 $F$ 为复测度 $\mu \in M[-\pi, \pi]$ 的 Cauchy--Stieltjes 积分,即
\begin{equation*}
F(z) = \frac{1}{2\pi}\int_{-\pi}^{\pi} C(r, \theta - t)\dd \mu(t) = \frac{1}{2\pi}\int_{-\pi}^{\pi} \frac{e^{it}}{e^{it} - z}\dd \mu(t),
\end{equation*}
则 $F(z)$ 在 $|z| \neq 1$ 时正则.

\begin{theorem}\label{thm:2.11.3}
若 $F$ 为复测度 $\mu \in M[-\pi, \pi]$ 的 Cauchy--Stieltjes 积分,则 $F \in \bigcap_{p<1} H^{p}$.
\end{theorem}

\begin{proof}
可设 $\mu$ 为正测度,此时由命题~\ref{prop:2.2.2} 有 $\operatorname{Re} F(z) > 0$.由定理~\ref{thm:2.11.2},$F \in \bigcap_{p<1} H^{p}$.
\end{proof}

已知在定理~\ref{thm:2.11.3} 中,其逆不成立,而且也不能把 $\bigcap_{p<1} H^{p}$ 换成 $H^{1}$.
